\documentclass[10pt,a4paper]{amsart}
\usepackage[margin=19mm]{geometry}
\usepackage{amsmath,amssymb,mathtools}
\usepackage[scr=rsfs,cal=euler]{mathalfa}
\usepackage{xfrac}
\usepackage[unicode,hidelinks,bookmarks=false]{hyperref}
\newcommand{\ie}{i.e.\ }
\newcommand{\cf}{cf.\ }
\newcommand{\resp}{respectively\ }
\newcommand{\Subsection}{\subsection}
\providecommand{\nobreakdash}{\nobreak\hbox{-}}
\setlength{\emergencystretch}{2em}
\multlinegap0pt
\usepackage{tikz,xcolor,aliascnt}
\usetikzlibrary{calc,through,backgrounds,shapes,matrix,math,arrows,trees}
\usepackage[nameinlink]{cleveref}
\newtheorem{theorem}{Theorem}[section]
\newaliascnt{lemma}{theorem}
\newtheorem{lemma}[lemma]{Lemma}
\aliascntresetthe{lemma}
\crefname{lemma}{Lemma}{Lemmas}
\newaliascnt{prop}{theorem}
\newtheorem{prop}[prop]{Proposition}
\aliascntresetthe{prop}
\crefname{prop}{Proposition}{Propositions}
\newaliascnt{coro}{theorem}
\newtheorem{coro}[coro]{Corollary}
\aliascntresetthe{coro}
\crefname{coro}{Corollary}{Corollarys}
\newaliascnt{defi}{theorem}
\newtheorem{defi}[defi]{Definition}
\aliascntresetthe{defi}
\crefname{defi}{Definition}{Definitions}
\newaliascnt{rema}{theorem}
\newtheorem{rema}[rema]{Remark}
\aliascntresetthe{rema}
\crefname{rema}{Remark}{Remarks}
\newenvironment{enonce*}[2][remark]{\par\medskip\noindent\textit{#2.}\ }{\par}
\definecolor{darkblue}{rgb}{0.0,0,0.7} 
\newcommand{\darkblue}{\color{darkblue}} 

\newcommand{\defn}[1]{\emph{\darkblue #1}}

\newcommand{\pop}{\mathsf{Pop}}
\newcommand{\wo}{w_\circ}
\newcommand{\U}{{W'}}
\newcommand{\uu}{{w'}}

\DeclareMathOperator{\spn}{span}
\DeclareMathOperator{\NC}{NC}
\DeclareMathOperator{\Des}{Des}
\DeclareMathOperator{\mov}{Mov}
\DeclareMathOperator{\fold}{fold}
\DeclareMathOperator{\im}{im}
\DeclareMathOperator{\unfold}{unfold}
\DeclareMathOperator{\cyc}{cyc}
\DeclareMathOperator{\aexc}{aexc}
\DeclareMathOperator{\Aexc}{Aexc}
\DeclareMathOperator{\Sym}{Sym}
\DeclareMathOperator{\Weak}{Weak}


\DeclarePairedDelimiter\abs{\lvert}{\rvert}

\newcommand*{\mk}{\mkern -1mu}
\newcommand*{\Mk}{\mkern -2mu}
\newcommand*{\mK}{\mkern 1mu}
\newcommand*{\MK}{\mkern 2mu}

\hypersetup{urlcolor=purple, linkcolor=blue, citecolor=red}


\newcommand*{\romanenumi}{\renewcommand*{\theenumi}{\roman{enumi}}}
\newcommand*{\Romanenumi}{\renewcommand*{\theenumi}{\Roman{enumi}}}
\newcommand*{\alphenumi}{\renewcommand*{\theenumi}{\alph{enumi}}}
\newcommand*{\Alphenumi}{\renewcommand*{\theenumi}{\Alph{enumi}}}

\title[Type-D pop-tsack termination]{Coxeter Pop-Tsack Torsing: original Theorem1.5}
\author{Colin Defant and Nathan Williams}
\date{Source: Algebraic Combinatorics5(3)(2022),559--581; DOI10.5802/alco.226}
\begin{document}\maketitle
\noindent\textit{Editorial scope.} The counted unit is original Theorem1.5, the bound for every type-D orbit. The universal inverse-element orbit is uncounted supporting core already present as a separate unit; it is retained here because it gives the sharp lower bound. The complete circle/central-label definitions and five-case author proof are retained. Figures consisting only of n=7 examples are omitted because all positions, blocks and cyclic orders used by the proof are fully specified in the original text. No original mathematical expression or hypothesis is repaired.
\section{Original operator and results}
\setcounter{equation}{2}
Let $T$ be the set of all reflections of $W$, and fix a Coxeter element $c$ of $W$ (a regular element of order $h$).  Let $\leq_T$ denote the absolute order on $W$. There is a general ``dual'' philosophy for finite Coxeter groups~\cite{bessis2003dual} in which the set $S$ of simple reflections is replaced by the set $T$ of all reflections, the longest element $\wo$ is replaced by the Coxeter element $c$, and the (left) weak order is replaced by the noncrossing partition lattice $\NC(W,c)$ (the interval $[e,c]$ in absolute order).

Define the \defn{noncrossing projection}
\begin{align}\label{EqNCProj}
\pi_T(\cdot,c)&:W \to \NC(W,c) \nonumber\\
\pi_T(w,c) &= \bigvee_{t \leq_T w}^{\NC(W,c)} t.
\end{align}

\setcounter{theorem}{1}
\begin{defi}
\label{def:dual_pop_stack}
The \defn{Coxeter pop-tsack torsing operator}\footnote{We have interchanged the roles of $S$ and $T$.} is the map
\begin{gather*}
 \pop_T(\cdot,c):W\to W \nonumber \\
\pop_T(w,c) = w \cdot \pi_T(w,c)^{-1}.
\end{gather*}
\end{defi}

When $c$ is understood, we may abbreviate $\pi_T(w,c)$ and $\pop_T(w,c)$ as $\pi_T(w)$ and $\pop_T(w)$, respectively; we establish in~\Cref{sec:general_properties} many structural properties of $\pop_T$, including that its orbit structure does not depend on $c$.  
By construction, the elements $w\in W$ such that $\pop_T(w)=e$ are exactly the noncrossing partitions---hence, the number of such elements is the \defn{$W$-Catalan number} 
\[
\left|\NC(W,c)\right|=\prod_{i=1}^n \frac{h+d_i}{d_i},
\]
 where $d_1,\ldots,d_n$ are the degrees of $W$.

\begin{theorem}
\label{thm:cinverseorbit}
The forward orbit of $c^{-1}$ under $\pop_T(\cdot,c)$ is 
\[
O_{\pop_T}(c^{-1})=\{c^{-1}, c^{-2},\ldots, c^{-(h-1)}, e\}.
\]
 The element $c^{-1}$ has no preimages under $\pop_T$. Moreover, if $2\leq i\leq h-1$, then the only preimage of $c^{-i}$ under $\pop_T$ is $c^{-(i-1)}$. 
\end{theorem}

\setcounter{theorem}{4}
\begin{theorem}
\label{thm:main2}
Let $W$ be a Coxeter group of type $D_n$. Then 
\begin{itemize}
    \item $\pop_T^{2n-3}(w)=e$ for all $w\in W$, and
    \item  $\max\limits _{w\in W}\left|O_{\pop_T}(w)\right|=2n-2$.
\end{itemize}
\end{theorem}
%\addtocounter{cdrthm}{-1}
%}

\section{Preliminaries}
Throughout this article, we assume that $W$ is a finite irreducible Coxeter group with a set $S$ of simple reflections. The \defn{rank} of $W$ is $|S|$. Given $g,w\in W$, we write $w^g$ for the conjugate $gwg^{-1}$ of $w$ by $g$. Let $T=\{s^w:s\in S, w\in W\}$ denote the set of \defn{reflections} of $W$. A \defn{$T$-word} for an element $w\in W$ is a word $t_1\cdots t_r$ over the alphabet $T$ that is equal to $w$ when considered as an element of $W$. The \defn{reflection length} of $w$, denoted $\ell_T(w)$, is the minimum length of a $T$-word for $w$. The \defn{absolute order} on $W$, denoted $\leq_T$, is defined by saying $v\leq_T w$ if $\ell_T(wv^{-1})=\ell_T(w)-\ell_T(v)$. Since $T$ is closed under conjugation by elements of $W$, we have $v\leq_T w$ if and only if $\ell_T(v^{-1}w)=\ell_T(w)-\ell_T(v)$. 

It is often useful to consider the \defn{geometric representation} of $W$ on a vector space $V\cong\mathbb R^{|S|}$. We refer the reader to \cite{armstrong2009generalized, BjornerBrenti} for more thorough discussions of this representation; here, we simply mention the basic results that we will need. We denote the action of an element $w\in W$ on a vector $\alpha\in V$ by $w\cdot\alpha$. For every $t\in T$, there is a hyperplane $H_t\subseteq V$ such that $t$ acts via the linear transformation that reflects through $H_t$. Furthermore, $w\cdot H_t=H_{t^w}$ for every $w\in W$ and $t\in T$. Define the \defn{moved space} of $w$ to be $\mov(w)=\im(w-1)$, the image of the linear transformation $w-1:V\to V$ ($1$ is the identity transformation). Brady and Watt \cite{BradyWatt} showed that if $t\in T$ and $w\in W$, then \begin{equation}\label{EqBradyWatt}
\mov(t)\subseteq\mov(w)\text{ if and only if }t\leq _T w.
\end{equation}

The ring $\Sym(V_\mathbb{C}^*)^W$ of $W$-invariant polynomials is a polynomial ring $\mathbb{C}[f_1,\ldots,f_n]$ over a set of invariant polynomials $f_i$ that may be chosen homogeneous and whose degrees $d_1 \leq \cdots \leq d_n$ are uniquely determined.  We write $h=d_n$ for the \defn{Coxeter number} of $W$.

\looseness-1
A \defn{standard Coxeter element} of $W$ is an element obtained by multiplying the simple reflections together in some order.  A \defn{regular element} of $W$ is an element with an eigenvector that lies in the complement of the reflecting hyperplanes.  Following~\cite{reiner2017non}, we define a \defn{Coxeter element} of $W$ to be a regular element of order $h$.  For Weyl groups, Coxeter elements coincide with conjugates of standard Coxeter elements.  
More generally, a \defn{reflection automorphism} of $W$ is a group automorphism that preserves the set of reflections of $W$.  By~\cite[Proposition~1.4]{reiner2017non}, the reflection automorphisms of $W$ act transitively on its Coxeter elements---thus, all Coxeter elements have the same order $h$.

We assume throughout this article that we have fixed a Coxeter element $c$ of $W$. A \defn{noncrossing partition} is an element $v\in W$ such that $v\leq_T c$. The interval $[e,c]$ in the absolute order is called the \defn{noncrossing partition lattice} and denoted $\NC(W,c)$; it is indeed a lattice. This is the only type
of lattice we will consider in the rest of the article, so all meets (denoted by $\bigwedge$) and joins (denoted by $\bigvee$) will be with respect to this lattice. Furthermore, the definitions of the noncrossing projection $\pi_T$ and the Coxeter pop-tsack torsing operator $\pop_T$ from \eqref{EqNCProj} and \Cref{def:dual_pop_stack} are made with respect to this 
fixed Coxeter element. Since the reflection automorphisms of $W$ act transitively on the Coxeter elements, different Coxeter elements give rise to isomorphic noncrossing partition lattices.

\section{Dihedral Groups}
\label{sec:dihedral}

As the dihedral case will be useful when analyzing the general case, we immediately dispense with it in this short section.  The dihedral group $I_2(h)$ has $2h$ elements: the identity $e$, $h$ reflections, a Coxeter element $c$, and then the elements $c^{-1},c^{-2},\ldots,c^{-(h-1)}$.  The identity, reflections, and Coxeter element are all noncrossing partitions, and they all map to the identity under a single application of $\pop_T$.  For each $1\leq i\leq h-1$, the element $c^{-i}$ requires exactly $h-i$ iterations of $\pop_T$ to reach the identity.

\setcounter{figure}{4}
\section{General Properties of \texorpdfstring{$\pop_T$}{Pop}}
\label{sec:general_properties}

\subsection{Absolute monotonicity}

Given $x\in W$, let $W_x$ be the subgroup of $W$ generated by the reflections lying weakly below $x$ in the absolute order. The subgroups $W_x$ for $x\in \NC(W,c)$ are called the \defn{noncrossing parabolic subgroups} of $W$. As discussed in \cite{ReadingNoncrossing}, the noncrossing parabolic subgroups of $W$ form a lattice under the containment order, and the map $x\mapsto W_x$ gives an isomorphism from $\NC(W,c)$ to this lattice. 

\begin{lemma}
\label{lem:monotone}
For every $x \in W$, we have $\pi_T(\pop_T(x))\leq_T \pi_T(x)$.
\end{lemma}
\begin{proof}
The parabolic subgroup generated by the reflections lying weakly below $\pi_T(x)$ in absolute order is $W_{\pi_T(x)}$. All of the reflections lying weakly below $x$ also lie weakly below $\pi_T(x)$, so $x$ is an element of $W_{\pi_T(x)}$. Of course, $\pi_T(x)$ also belongs to $W_{\pi_T(x)}$. This means that $\pop_T(x)=x\pi_T(x)^{-1}$ is in $W_{\pi_T(x)}$, so all of the reflections weakly below $\pop_T(x)$ are in $W_{\pi_T(x)}$. It follows that 
$W_{\pi_T(\pop_T(x))}$ is contained in $W_{\pi_T(x)}$, and this proves the result. 
\end{proof}

The following proposition gives a condition for elements to have an ``isolated'' forward orbit under $\pop_T$. It states that the only possible preimage under $\pop_T$ of an element $w \in W$ with $\pi_T(w)=c$ is $wc$.

\begin{prop}
\label{prop:preimages}
If $x,w\in W$ are such that $\pop_T(x)=w$ and $\pi_T(w)=c$, then $x=wc$.
\end{prop}

\begin{proof}
Suppose $x$ is such that $\pop_T(x)=w$.  By \Cref{lem:monotone}, $c\leq_T\pi_T(x)$, so we must have $\pi_T(x)=c$. However, this implies that $x=\pop_T(x)\pi_T(x)=wc$.  
\end{proof}

\subsection{Equivariance under reflection automorphisms}

In this section, we make a point of decorating $\pi_T$ and $\pop_T$ by the Coxeter element $c$ in order to differentiate different choices of Coxeter elements.  We show a sort of equivariance of $\pop_T$ with respect to reflection automorphisms, which establishes that the orbit structure of $\pop_T$ does not depend on the choice of Coxeter element $c$.  Thus, the Coxeter pop-tsack torsing operators defined with respect to different Coxeter elements are dynamically equivalent to each other, so none of our results about $\pop_T$ depend on the choice of $c$. This means that we have the freedom to choose any Coxeter element we want when we work with specific Coxeter groups, and it will often be helpful to use a specific choice that makes the resulting combinatorics easiest to describe.

\begin{prop}  Let $\phi$ be a reflection automorphism of $W$.
For $w \in W$, we have
\begin{itemize}
\item $\phi(\pi_T(w,\phi^{-1}(c))) = \pi_T(\phi(w),c)$, and
\item $\phi(\pop_T(w,\phi^{-1}(c))) = \pop_T(\phi(w),c)$.
\end{itemize}
\label{prop:pop_props}
\end{prop}

\begin{proof}
For the first point, observe that the map $w \mapsto \phi(w)$ is a lattice isomorphism from $\NC(W,c)$ to $\NC(W,\phi(c))$.  Furthermore, for each reflection $t$, we have $t \leq_T w$ if and only if $\phi(t) \leq_T \phi(w)$.  Therefore, 
\[
\phi\left(\pi_T(w,\phi^{-1}(c))\right) = \phi\left(\bigvee_{t \leq_T w}^{\NC(W,\phi^{-1}(c))} t\right) = \bigvee_{t \leq_T w}^{\NC(W,c)} \phi(t) = \bigvee_{t \leq_T \phi(w)}^{\NC(W,c)} t = \pi_T(\phi(w),c).\]  


For the second point, we use the first point to check that 
\begin{align*}\pop_T(\phi(w),c) &= \phi(w) \pi_T(\phi(w),c)^{-1} = \phi(w) \left(\phi\left(\pi_T(w,\phi^{-1}(c))\right)\right)^{-1}\\ & = \phi\left(w\pi_T(w,\phi^{-1}( c))^{-1}\right) = \phi\left(\pop_T(w,\phi^{-1}(c))\right) .\qedhere
\end{align*}
\end{proof}

As conjugation by $c$ is a special kind of reflection automorphism that restricts to a lattice isomorphism of $\NC(W,c)$, we can specialize the preceding result to conclude the following.

\begin{coro} For $w \in W$, we have
\begin{itemize}
    \item  $\pi_T(w,c)^c = \pi_T(w^c,c)$, and
    \item $\pop_T(w,c)^c = \pop_T(w^c,c)$. 
\end{itemize}
\end{coro}


Since reflection automorphisms act transitively on the set of Coxeter elements, we further obtain:
\begin{coro}
\label{cor:orbit_structure}
The orbit structure of $\pop_T$ does not depend on the choice of Coxeter element.
\end{coro}

\subsection{Folding and the Coxeter plane}
\label{sec:folding}

 


\begin{defi}
Let $(W,S)$ and $(\U,S')$ be Coxeter systems such that $W$ and $W'$ are finite and irreducible. Suppose there is a surjective map $\fold:S\to S'$ such that for each $s'\in S'$, the elements of $\fold^{-1}(s')$ all commute with each other. If there is an injective homomorphism $\unfold:\U\to W$ with $\unfold(s')=\prod_{s\in\fold^{-1}(s')}s$ for all $s'\in S'$, then we say that $W$ \defn{folds} to $W'$ and that $W'$ \defn{unfolds} into $W$. 
\end{defi}

We remark that our definition of folding is more general than the standard definition that uses an automorphism of the Dynkin diagram. 

\Cref{fig:folding} illustrates several examples of folding. If $W$ folds to $\U$ and $c'$ is a Coxeter element of $\U$, then $\unfold(c')$ is a Coxeter element of $W$. Hence, the Coxeter number is preserved by folding.


A useful example of folding goes through the Coxeter plane.  Any finite irreducible Coxeter group $W$ with Coxeter number $h$ folds to the dihedral group $I_2(h)$ of order $2h$, as we now explain.  Since the Coxeter--Dynkin diagram of $W$ is a tree (hence, a bipartite graph), the set $S$ of simple reflections of $W$ naturally separates into two sets $S_+$ and $S_-$, where the simple reflections within each set all commute with each other.  Writing $c_\pm$ for the product of the reflections in $S_\pm$ so that $c = c_+ c_-$ is a Coxeter element of $W$, $\langle c_+,c_-\rangle$ is the dihedral group $I_2(h)$.  Each of $c_+$ and $c_-$ acts as a reflection on the \defn{Coxeter plane}, a plane spanned by the real and imaginary parts of an eigenvector for $c$ with eigenvalue $e^{2\pi i/h}$.


\begin{figure}[htbp]
    \centering
    \bgroup
\def\arraystretch{2}
\begin{tabular}{|ccc|c|}
\hline
     $W$ & $\U$ & $h$ & Diagrammatic folding  \\ \hline
     $A_{2n-1}$ & $B_n$ & $2n$\vrule height20pt depth20pt width0pt & \raisebox{-0.5\height}{\begin{tikzpicture}[scale=.5]
       \filldraw[black] (0,.5) circle (2pt) -- (1,0) circle (2pt) -- (2,0) circle (2pt);
        \filldraw[black,dotted] (2,0) circle (2pt) -- (3,0) circle (2pt);
        \filldraw[black] (3,0) circle (2pt) -- (4,0) circle (2pt);
        \filldraw[black] (0,.5) circle (2pt) -- (1,1) circle (2pt) -- (2,1) circle (2pt);
        \filldraw[black,dotted] (2,1) circle (2pt) -- (3,1) circle (2pt);
        \filldraw[black] (3,1) circle (2pt) -- (4,1) circle (2pt);
        \filldraw[black,double] (0,-.5) -- (1,-.5);
        \filldraw[black] (0,-.5) circle (2pt);
        \filldraw[black] (1,-.5) circle (2pt);
        \filldraw[black] (1,-.5) circle (2pt) -- (2,-.5) circle (2pt);
        \filldraw[black,dotted] (2,-.5) circle (2pt) -- (3,-.5) circle (2pt);
        \filldraw[black] (3,-.5) circle (2pt) -- (4,-.5) circle (2pt);
      \end{tikzpicture}}\\ 
     $D_n$ & $B_{n-1}$ & $2n-2$\vrule height0pt depth20pt width0pt & \raisebox{-0.5\height}{\begin{tikzpicture}[scale=.5]
     \filldraw[black] (0,1) circle (2pt) -- (1,.5) circle (2pt);
       \filldraw[black] (0,0) circle (2pt) -- (1,.5) circle (2pt) -- (2,.5) circle (2pt);
        \filldraw[black,dotted] (2,.5) circle (2pt) -- (3,.5) circle (2pt);
        \filldraw[black] (3,.5) circle (2pt) -- (4,.5) circle (2pt);
        \filldraw[black,double] (0,-.5) -- (1,-.5);
        \filldraw[black] (0,-.5) circle (2pt);
        \filldraw[black] (1,-.5) circle (2pt);
        \filldraw[black] (1,-.5) circle (2pt) -- (2,-.5) circle (2pt);
        \filldraw[black,dotted] (2,-.5) circle (2pt) -- (3,-.5) circle (2pt);
        \filldraw[black] (3,-.5) circle (2pt) -- (4,-.5) circle (2pt);
      \end{tikzpicture}} \\ 
     $E_6$ & $F_4$ & $12$\vrule height0pt depth25pt width0pt & \raisebox{-0.5\height}{\begin{tikzpicture}[scale=.5]
       \filldraw[black] (-1,.5) circle(2pt) -- (0,.5) circle (2pt) -- (1,0) circle (2pt) -- (2,0) circle (2pt);
        \filldraw[black] (0,.5) circle (2pt) -- (1,1) circle (2pt) -- (2,1) circle (2pt);
        \filldraw[black,double] (0,-.5) -- (1,-.5);
        \filldraw[black] (-1,-.5) circle (2pt) -- (0,-.5) circle (2pt);
        \filldraw[black] (0,-.5) circle (2pt);
        \filldraw[black] (1,-.5) circle (2pt);
        \filldraw[black] (1,-.5) circle (2pt) -- (2,-.5) circle (2pt);
      \end{tikzpicture}}\\ 
     $E_8$ & $H_4$ & 30\vrule height0pt depth20pt width0pt & \raisebox{-0.5\height}{\begin{tikzpicture}[scale=.5]
       \filldraw[black] (0,1.5) circle(2pt) -- (1,1) circle (2pt) -- (2,1) circle (2pt) -- (3,1) circle (2pt);
        \filldraw[black] (0,.5) circle (2pt) -- (1,0) circle (2pt) -- (2,0) circle (2pt) -- (3,0) circle (2pt);
        \filldraw[black] (0,.5) -- (1,1);
        \filldraw[black] (0,-.5) circle (2pt) -- (1,-.5) circle (2pt) -- (2,-.5) circle (2pt) -- (3,-.5) circle (2pt);
        \node at (.5,-1) {\scriptsize 5};
      \end{tikzpicture}}\\
      \hline
\end{tabular}
\egroup
    \caption{Examples of folding of finite Coxeter groups.}
    \label{fig:folding}
\end{figure}

In this subsection, we write $T$ and $T'$ for the sets of reflections in $W$ and $W'$, respectively. 

\begin{lemma}
Let $\U$ be a finite irreducible Coxeter group that unfolds into $W$. Let $c'$ be a bipartite Coxeter element of $\U$, and let $c=\unfold(c')$ be the corresponding Coxeter element of $W$. Then $\unfold(\NC(\U,c'))$ is a sublattice of $\NC(W,c)$.
\label{lem:lattice_embedding}
\end{lemma}
\begin{proof}
Suppose that $w'\in W'$ and $w' \leq_{T'} c'$ so that $\uu$ can be found as a prefix of $c'$. Then $\unfold(\uu)$ can be found as a prefix of $c$, so $\unfold(\uu) \leq_T c$.  By the same reasoning (replacing now $c'$ by any element of $\NC(\U,c')$), we have that $\unfold(\NC(\U,c'))$ is a subposet of $\NC(W,c)$.

We now check that $\unfold(\NC(\U,c'))$ is a sublattice of $\NC(W,c)$ (that is, that $\unfold(\NC(\U,c'))$ is closed under the meet and join operations of $\NC(W,c)$).  We know by~\cite{brady2008non} that a noncrossing partition $w$ is characterized by the set of reflections that lie below it: $w = \bigvee_{t \leq_T w} t$. Since 
\[
w \wedge v = \bigvee\{t : t \leq_T w \text{ and } t \leq_T v\},
\]
 it follows that $\unfold(\NC(\U))$ is a meet-sublattice. Indeed, for fixed $t'\in T'$ and $r \in \fold^{-1}(t')$, we have $r \leq \unfold(w') \wedge \unfold(v')$ if and only if $r \leq_T \unfold(w')$ and $r \leq_T \unfold(v')$. This occurs if and only if $t' \leq_{T'} w'$ and $t' \leq_{T'} v'$, which occurs if and only if $t' \leq_{T'} w' \wedge v'$, which occurs if and only if $r \leq_T \unfold(w' \wedge v')$. 


The Kreweras complement maps $K:\NC(W,c)\to \NC(W,c)$ and $K':\NC(W',c')\to \NC(W',c')$, defined by $K(w)= cw^{-1}$ and $K'(w')=c'(w')^{-1}$, are lattice anti-isomorphisms \cite{armstrong2009generalized}. We have $K\circ\unfold=\unfold\circ K'$ because $\unfold$ is a homomorphism that sends $c'$ to $c$. This shows that $K$ preserves the set $\unfold(\NC(\U,c'))$, so $\unfold(\NC(\U,c'))$ is also a join-sublattice of $\NC(W,c)$.
\end{proof}

\setcounter{subsection}{3}
\subsection{The inverse Coxeter element orbit}
\label{sec:inverse_c_orbit}

\setcounter{theorem}{11}
\begin{lemma}\label{Lem:joinofcorbit}
If $t\in T$ is a reflection, then $c=\bigvee\limits_{k=0}^{h-1} t^{c^{k}}$.
\end{lemma}

\begin{proof}
Since $W$ folds to $I_2(h)$ via the Coxeter plane, we can consider the Coxeter element $\unfold(c')$ of $W$, where $c'$ is a Coxeter element of $I_2(h)$. Each Coxeter element of $W$ can be obtained by applying a reflection automorphism to $\unfold(c')$, so it suffices to prove the lemma when $c=\unfold(c')$.

Conjugation by $c$ induces an automorphism of $\NC(W,c)$; therefore, if a set $O \subseteq \NC(W,c)$ is invariant under conjugation by $c$, then the element $x=\bigvee_{w \in O} w$ is invariant under conjugation by $c$:
	\[
x^c = \left(\bigvee_{w \in O} w\right)^c = \bigvee_{w \in O} w^c = \bigvee_{w \in O} w = x.
\]
It is a classical result that the centralizer of $c$ is the cyclic group generated by $c$ (see \cite[Proposition~30]{Carter}), so we must have $x=c^i$ for some integer $i$.  Therefore, we are left to prove that the only powers of $c$ in $\NC(W,c)$ are $c$ and $e$. This will follow if we can show that the only powers of $c'$ in $\NC(I_2(h),c')$ are $c'$ and $e$.  But this is clear since the absolute length of any power of $c'$ (a rotation) is $0$ or $2$ (see also~\cite[Lemma~12.2]{bessis2015finite}).
\end{proof}

\begin{prop}
For each integer $1\leq i\leq h-1$, we have $\pop_T(c^{-i})=c^{-(i+1)}$. Thus, the forward orbit of $c^{-1}$ under $\pop_T$ is $\{c^{-1},c^{-2},\ldots,c^{-(h-1)},e\}$, which has size~$h$.
\label{prop:forward_cinv}
\end{prop}

\begin{proof}
Fix $1\leq i\leq h-1$. We claim that if $t\in T$ is such that $t \leq_T c^{-i}$, then $t^c \leq_T c^{-i}$.  According to \eqref{EqBradyWatt}, we have $\mov(t) \subseteq \mov(c^{-i})$.  Take $\alpha_t$ to be a vector orthogonal to the hyperplane $H_t$ so that $\mov(t) = \spn\{\alpha_t\}$.  Then $\alpha_t\in\mov(c^{-i})$, so there exists a vector $\beta\in V$ such that $(c^{-i}-1)\cdot\beta = \alpha_t$. Applying $c$ to each side of this equation shows that 
\[
(c^{-i}-1)\cdot(c\cdot\beta)=c\cdot((c^{-i}-1)\cdot\beta)=c\cdot\alpha_t.
\]
 Thus, $\spn\{c\cdot\alpha_t\}\subseteq\mov(c^{-i})$. Since $c\cdot H_t=H_{t^c}$, we know that $c\cdot\alpha_t$ is orthogonal to $H_{t^c}$. Thus, $\mov(t^c)=\spn\{c\cdot\alpha_t\}$. Appealing to \eqref{EqBradyWatt} again, we find that $t^c\leq_T c^{-i}$. 

Note that there must exist some $t\in T$ with $t \leq_T c^{-i}$ since $c^{-i} \neq e$. It follows from the preceding paragraph that $t^{c^k}\leq_T c^{-i}$ for all $0\leq k\leq h-1$. Therefore, we can use \Cref{Lem:joinofcorbit} to see that 
\[
c=\bigvee\limits_{k=0}^{h-1} t^{c^{k}}\leq_T\bigvee_{t'\leq_T c^{-i}}t'=\pi_T(c^{-i}).
\]
 Consequently, $\pi_T(c^{-i})=c$, so $\pop_T(c^{-i})=c^{-i}c^{-1}=c^{-(i+1)}$. 
\end{proof}

\noindent\textit{Original closing inverse-element argument.}
\begin{proof}
\looseness-1
  By \Cref{prop:preimages} and \Cref{prop:forward_cinv}, the only possible preimage of $c^{-i}$ under $\pop_T$ is $c^{-(i-1)}$.  Furthermore, $c^{-(i-1)}$ is actually a preimage of $c^{-i}$ unless $i=1$.
\end{proof}

This completes the proof of \Cref{thm:cinverseorbit}.

\setcounter{section}{5}
\section{Type D}
\label{sec:typed}
In this section, we prove the following theorem.

\setcounter{theorem}{1}
As in the previous section, let $\pm[n]=\{\overline n,\ldots,\overline 1,1,\ldots,n\}$ (with bars replacing negative signs). The Coxeter group $D_n$ is the group of permutations $w:\pm[n]\to\pm[n]$ such that $w(\overline i)=\overline{w(i)}$ for all $i\in\pm[n]$ and such that $\abs{\{i\in[n]:w(i)<0\}}$ is even. The reflections in $D_n$ are the elements of the form $(i\,\,j)(\overline i\,\,\overline j)$, where $i,j\in\pm[n]$ and $i\neq\overline j$. Throughout this section, we fix the Coxeter element 
$c=(\overline 1\,\,\overline 2\,\cdots\overline{n-1}\,\,1\,\,2\cdots n-1)(\overline n\,\,n)$. Given a cycle $\mathcal C=(i_1i_2\cdots i_k)$ in an element $w\in D_n$, let $\overline{\mathcal C}$ denote the cycle $(\overline{i_1}\,\,\overline{i_2}\cdots \overline{i_k})$, and note that $\overline{\mathcal C}$ is also a cycle in $w$. We say the cycle $\mathcal C$ is \defn{balanced} if $\mathcal C=\overline{\mathcal C}$. It follows from the definition of $D_n$ that $\mathcal C$ is balanced if and only if there exists an entry $k$ such that $k$ and $\overline k$ are both in $\mathcal C$. A cycle that is not balanced is \defn{unbalanced}. Each element of $D_n$ has an even number of balanced cycles. 

Every element $w\in D_n$ can be written uniquely as 
\[
w=\mathcal C_1\overline{\mathcal C_1}\cdots\mathcal C_r\overline{\mathcal C_r}\mathcal D_1\cdots\mathcal D_s,
\]
 where $\mathcal C_1,\ldots,\mathcal C_r$ are unbalanced cycles, $\mathcal D_1,\ldots,\mathcal D_s$ are balanced cycles, and the cycles appearing in the factorization are pairwise disjoint. The reflection length of $w$ can be computed as 
\[
\ell_T(w)=\sum_{i=1}^r(k_i-1)+\sum_{j=1}^sm_j,
\]
 where $k_i$ is the number of entries in $\mathcal C_i$ (and, therefore, also $\overline{\mathcal C_i}$) and $m_j$ is the number of entries in $\mathcal D_j$ (see \cite[Section~3]{BradyWatt}). The following useful lemma is a consequence of the computations in \cite[Example~3.6]{BradyWatt} (see also \cite[Section~3]{AthanasiadisReiner}).

\begin{lemma}[\cite{BradyWatt}]\label{Lem1}
Let $w\in D_n$, and consider $i,j\in\pm[n]$ with $i\not\in\{\overline j,j\}$. We have $(i\,\,j)(\overline i\,\,\overline j)\leq_T w$ if and only if $i$ and $j$ are in the same cycle of $w$ or $i$ and $j$ are in different balanced cycles of $w$. 
\end{lemma}

Our combinatorial description of noncrossing partitions in type~$D$ is based on the work of Athanasiadis and Reiner \cite{AthanasiadisReiner}. Begin by labeling $2n{-}2$ equally-spaced points on a circle with the numbers $\overline 1,\overline 2,\ldots,$ $\overline{n-1},1,2,\ldots,n-1$ in clockwise order. Now place a single extra point at the center of the circle, and label it with both the number $n$ and the number $\overline n$. Consider a set partition $\rho$ of $\pm[n]$; we associate each block of $\rho$ with the set of points labeled by the elements of that block. We say $\rho$ is a \defn{type-$D_n$ noncrossing set partition} if it satisfies the following conditions: 


\begin{enumerate}
    \item\label{sect5_1} For every block $B$ of $\rho$, the set $\overline B=\{\overline i:i\in B\}$ is also a block of $\rho$.
    \item\label{sect5_2} If $B$ and $B'$ are distinct blocks of $\rho$, then the relative interior of the convex hull of $B$ does not contain a point of the convex hull of $B'$. 
    \item\label{sect5_3} The set $\{\overline n,n\}$ is not a block of $\rho$. 
\end{enumerate}

We associate each type-$D_n$ noncrossing set partition $\rho$ with the diagram showing the labeled points on the circle and at the center of the circle together with the convex hulls of the blocks of $\rho$. A \defn{zero block} of $\rho$ is a block $Z$ of $\rho$ such that $Z=\overline Z$. Note that $\rho$ has at most $1$ zero block; moreover, if $\rho$ does have a zero block $Z$, then $Z$ must properly contain the set $\{\overline n,n\}$. 

Suppose $\rho$ is a type-$D_n$ noncrossing set partition. For each block $B$ of $\rho$ that is not a zero block, form a cycle by reading the elements of $B$ in clockwise order around the boundary of the convex hull of $B$. If $\rho$ has a zero block $Z$, then form the cycle $(\overline n\,\, n)$, and form another cycle by reading the elements of $Z\setminus\{\overline n,n\}$ in clockwise order around the boundary of the convex hull of $Z$. Multiplying all of these cycles together forms the disjoint cycle decomposition of a permutation $g(\rho)$ of $\pm[n]$. Athanasiadis and Reiner \cite{AthanasiadisReiner} proved that $g(\rho)$ is actually a noncrossing partition in $\NC(D_n,c)$ and that the map $g$ provides a bijection from the set of type-$D_n$ noncrossing set partitions to the set $\NC(D_n,c)$. 


\begin{lemma}\label{Lem3}
Suppose $v\in \NC(D_n,c)$ and $k\in\pm[n-1]$ are such that $v(k)\in\{\overline n,n\}$. Let $Y$ be the set consisting of the next $n-1$ numbers after $k$ in clockwise order around the circle. Then none of the elements of $Y$ lie in the same cycle as $k$ in $v$.  
\end{lemma}

\begin{proof}
Let $\mathcal C$ be the cycle in $v$ containing $k$. We can write $v=g(\rho)$ for some type-$D_n$ noncrossing set partition $\rho$. The entries in $\mathcal C$ form a block $B$ of $\rho$ that is not a zero block. By the definition of $g$, the cycle $\mathcal C$ is formed by reading the elements of $B$ in clockwise order around the boundary of the convex hull of $B$. This immediately implies that $\mathcal C\cap Y=\emptyset$.  
\end{proof}

\begin{lemma}\label{Lem2}
Let $u\in D_n$, and let $v=\pi_T(u)$. Consider distinct $i,j\in\pm[n-1]$. If $i$ and $j$ are in the same cycle of $u$, then they are in the same cycle of $v$.   
\end{lemma}

\begin{proof}
Suppose first that $i\neq \overline j$. It follows from \Cref{Lem1} that $(i\,\,j)(\overline i\,\,\overline j)\leq_T u$, so $(i\,\,j)(\overline i\,\,\overline j)\leq_T v$. Since $v\in\NC(D_n,c)$ is noncrossing, it has at most $1$ balanced cycle other than $(\overline n\,\,n)$. This means that $i$ and $j$ cannot belong to different balanced cycles of $v$, so they must belong to the same cycle of $v$ by \Cref{Lem1}. 

Now suppose $i=\overline j$. This means $i$ is in a balanced cycle in $u$. We want to show that $i$ is in a balanced cycle of $v$. Suppose, by way of contradiction, that $i$ is in an unbalanced cycle of $v$. Each element of $D_n$ has an even number of balanced cycles, so there is some balanced cycle in $u$ that does not contain $i$; let $k$ be one of the numbers in this other balanced cycle. Since $i$ and $k$ are in different balanced cycles in $u$, we have $(i\,\,k)(\overline i\,\,\overline k)\leq_T u$ by \Cref{Lem1}. Therefore, $(i\,\,k)(\overline i\,\,\overline k)\leq_T v$. Because $i$ is in an unbalanced cycle in $v$, \Cref{Lem1} tells us that $i$ and $k$ are in the same cycle in $v$. Since $\overline k$ is in the same balanced cycle in $u$ as $k$, the same argument we have just given shows that $i$ and $\overline k$ are in the same cycle in $v$. However, this means that the cycle in $v$ containing $i$ contains both $k$ and $\overline k$, so it is balanced. This is a contradiction.     
\end{proof}

Our goal in the remainder of this section is to prove \Cref{thm:main2}. The Coxeter number of $D_n$ is $2n{-}2$. We already know by \Cref{prop:forward_cinv} that $O_{\pop_T}(c^{-1})$ is a forward orbit of $\pop_T:D_n\to D_n$ of size $2n{-}2$. Thus, we need to prove that $\pop_T^{2n-3}(u_0)=e$ for all $u_0\in D_n$. 

\looseness-1
Fix $u_0\in D_n$. For $k\geq 0$, let $u_k=\pop_T^k(u_0)$ and $v_k=\pi_T(u_k)$. Thus, $u_{k+1}=u_kv_k^{-1}$. Our goal is to show that $u_{2n-3}=e$. Since $\pop_T^{-1}(e)=\NC(D_n,c)$, it suffices to show that $u_{2n-4}\in\NC(D_n,c)$. The following lemma will do all of the heavy lifting. When we refer to ``the circle,'' we mean the circle that has $2n{-}2$ equally-spaced points labeled with the numbers $\overline 1,\overline 2,\ldots,\overline{n-1},1,2,\ldots,n-1$ in clockwise order and that has its center labeled with the numbers $\overline n$ and $n$. Define the \defn{predecessor} of an element $j\in\pm[n-1]$ to the be number that appears immediately before $j$ in the clockwise order of the circle. For example, the predecessor of $2$ is $1$, and the predecessor of $\overline 1$ is $n-1$.  

\begin{lemma}\label{lem:DMainLemma}
Suppose $j\in\pm[n-1]$ is such that $u_0^{-1}(j)\not\in\{\overline n,n\}$. Then $u_{2n-4}^{-1}(j)$ is either $j$ or the predecessor of $j$. 
\end{lemma}

\begin{proof}
Define a total order $\prec$ on $\pm[n-1]$ by $\overline 1\prec\cdots\prec\overline{n-1}\prec 1\prec\cdots\prec n-1$. Define $\prec_j$ to be the total order on $\pm[n-1]$ obtained by cyclically shifting the total order $\prec$ so that the maximum element in the order $\prec_j$ is $j$. In other words, $\prec_j$ is given by reading the numbers clockwise around the circle so that the last number read is $j$. Given a set $Q\subseteq\pm[n-1]$, we write $i\prec_j Q$ (respectively, $Q\prec_J i$) to mean that $i\prec_j q$ (respectively, $q\prec_j i$) for all $q\in Q$. 

If $u_k^{-1}(j)=j$ for some $0\leq k\leq 2n-4$, then $u_{k'}^{-1}(j)=j$ for all $k'\geq k$. In particular, $u_{2n-4}^{-1}(j)=j$ in this case. Thus, we may assume $u_k^{-1}(j)\neq j$ for all $0\leq k\leq 2n-4$.  

Suppose $k\in\{0,\ldots,2n-4\}$ is such that $u_k^{-1}(j)\not\in\{\overline n,n\}$ and $u_{k+1}^{-1}(j)=v_k(u_k^{-1}(j))\not\in\{\overline n,n\}$. Let $\mathcal C$ be the cycle in $v_k$ containing $j$. Since $u_k^{-1}(j)$ and $j$ are in the same cycle in $u_k$, \Cref{Lem2} tells us that $u_k^{-1}(j)\in\mathcal C$. Because $v_k$ is a noncrossing partition, the cycle $\mathcal C$ is ordered clockwise. Consequently, $u_{k+1}^{-1}(j)$ is the element of $\mathcal C$ that occurs right after $u_k^{-1}(j)$ in the clockwise cyclic order. Since $u_k^{-1}(j)\neq j$, we must have $u_k^{-1}(j)\prec_j u_{k+1}^{-1}(j)$.

If there does not exist an integer $i\in\{0,\ldots,2n-4\}$ such that $u_i^{-1}(j)\in\{\overline n,n\}$, then it follows from the preceding paragraph that $u_0^{-1}(j)\prec_j u_1^{-1}(j)\prec_j\cdots\prec_j u_{2n-4}^{-1}(j)$. We have assumed $u_{2n-4}^{-1}(j)\neq j$, so it follows that $u_{2n-4}^{-1}(j)$ is the predecessor of $j$, as desired. 

We may now assume there exists $i\in\{0,\ldots,2n-4\}$ with $u_i^{-1}(j)\in\{\overline n,n\}$. We will actually show that this case cannot occur because it contradicts our assumption that $u_{k}^{-1}(j)\neq j$ for all $0\leq k\leq 2n-4$. Since $u_0^{-1}(j)\not\in\{\overline n,n\}$ by the hypothesis of the lemma, we must have $i\geq 1$ and $u_{i-1}^{-1}(j)\not\in\{\overline n,n\}$. Let $\mathcal C_{i-1}$ and $\mathcal C_i$ be the cycles in $v_{i-1}$ and $v_i$, respectively, that contain $j$. 

\looseness-1
Let $Y$ be the set consisting of the next $n-1$ numbers after $u_{i-1}^{-1}(j)$ in clockwise order around the circle. Because $v_{i-1}(u_{i-1}^{-1}(j))=u_i^{-1}(j)\in\{\overline n,n\}$, we can use \Cref{Lem3} to see that none of the elements of $Y$ lie in the same cycle of $v_{i-1}$ as $u_{i-1}^{-1}(j)$. Since $j$ and $u_{i-1}^{-1}(j)$ are in the same cycle of $u_{i-1}$, we know by \Cref{Lem2} that $j$ and $u_{i-1}^{-1}(j)$ are in the same cycle of $v_{i-1}$. In other words, $u_{i-1}^{-1}(j)\in\mathcal C_{i-1}$, and $\mathcal C_{i-1}\cap Y=\emptyset$. Hence, $j\not\in Y$. This implies that $u_{i-1}^{-1}(j)\prec_j Y\prec_j j$. In particular, $u_{i-1}^{-1}(j)$ must be among the smallest $n-2$ numbers in the order $\prec_j$. We also know that $u_i^{-1}(j)\in\mathcal C_{i-1}$ because $u_i^{-1}(j)=v_{i-1}(u_{i-1}^{-1}(j))$. Since $u_i^{-1}(j)\in\{\overline n,n\}$, this implies that $\mathcal C_{i-1}$ is an unbalanced cycle in $v_{i-1}$. In fact, this shows that $v_{i-1}$ has no balanced cycles (if it did, then by the definition of the bijection $g$, one of the balanced cycles of $v_{i-1}$ would be $(\overline n\,\, n)$).  

In~\cite{AthanasiadisReiner}, Athanasiadis and Reiner gave an explicit combinatorial description of the covering relations in $\NC(D_n,c)$; it follows from their description that if $x,y\in\NC(D_n,c)$ are such that $x\leq_T y$, then every unbalanced cycle in $y$ is a union (as a set) of unbalanced cycles in $x$.

According to \Cref{lem:monotone}, we have $v_i\leq_T v_{i-1}$. Because $v_{i-1}$ has no balanced cycles, neither does $v_i$. Furthermore, $\mathcal C_i\subseteq \mathcal C_{i-1}$.  Since $j$ and $u_i^{-1}(j)$ are in the same cycle in $u_i$, we have $(j\,\,u_i^{-1}(j))(\overline j\,\,\overline{u_i^{-1}(j)})\leq_T u_i$ by \Cref{Lem1}. Therefore, $(j\,\,u_i^{-1}(j))(\overline j\,\,\overline{u_i^{-1}(j)})\leq_T v_i$. Since $v_i$ has no balanced cycles, we can use \Cref{Lem1} again to see that $j$ and $u_i^{-1}(j)$ are in the same unbalanced cycle in $v_i$. That is, $u_i^{-1}(j)\in\mathcal C_i$. Now, $u_{i+1}^{-1}(j)=v_i(u_i^{-1}(j))$ is the element of $\mathcal C_i$ that occurs immediately after $u_i^{-1}(j)$ in the clockwise cyclic order, so $Y\prec_j u_{i+1}^{-1}(j)$. Hence, $u_{i+1}^{-1}(j)$ is among the largest $n-2$ elements of $\pm[n-1]$ in the order $\prec_j$.  

We have shown that if $i\in\{1,\ldots,2n-4\}$ is such that $u_{i}^{-1}(j)\in\{\overline n,n\}$ and $u_{i-1}^{-1}(j)\not\in\{\overline n,n\}$, then $u_{i-1}^{-1}(j)$ is among the smallest $n-2$ elements of $\pm[n-1]$ in the order $\prec_j$ while $u_{i+1}^{-1}(j)$ is among the largest $n-2$ elements. We have also seen that if $k\in\{0,\ldots,2n-4\}$ is such that $u_k^{-1}(j)\not\in\{\overline n,n\}$, then $u_k^{-1}(j)\prec_j u_{k+1}^{-1}(j)$. We assumed in the hypothesis of the lemma that $u_0^{-1}(j)\not\in\{\overline n,n\}$. It follows from these facts that $u_k^{-1}(j)\not\in\{\overline n,n\}$ for every $k\in\{0,\ldots,2n-4\}\setminus\{i\}$. Therefore, the sequence $u_0^{-1}(j),u_1^{-1}(j),\ldots,u_{i-1}^{-1}(j),u_{i+1}^{-1}(j),\ldots,u_{2n-4}^{-1}(j)$ is strictly increasing in the order $\prec_j$. However, this sequence does not include any of the $n-1$ elements of $Y$ because $u_{i-1}^{-1}(j)\prec_j Y\prec_j u_{i+1}^{-1}(j)$. This is our desired contradiction. 
\end{proof}

\begin{proof}[Proof of \Cref{thm:main2}]
Earlier, we chose an arbitrary $u_0\in D_n$ and let $u_k=\pop_T^k(u_0)$. Our goal is to prove that $u_{2n-3}=e$; in order to do so, it suffices to show that $u_{2n-4}\in\NC(D_n,c)$. Let $m=u_0(n)$. \Cref{lem:DMainLemma} tells us that if $j\in\pm[n-1]\setminus\{\overline m,m\}$, then $u_{2n-4}^{-1}(j)$ is either $j$ or the predecessor of $j$. We now consider several cases. 




\begin{enonce*}[remark]{Case 1} Suppose $u_{2n-4}^{-1}(m)=m$. Then $u_{2n-4}^{-1}(\overline m)=\overline m$. If $m\not\in\{\overline n,n\}$, then one can straightforwardly deduce from \Cref{lem:DMainLemma} that $u_{2n-4}=e\in\NC(D_n,c)$. If $m\in\{\overline n,n\}$, then one can use \Cref{lem:DMainLemma} to see that $u_{2n-4}$ is either $e$ or $c$; in either case, $u_{2n-4}\in\NC(D_n,c)$. 
\end{enonce*}

\begin{enonce*}[remark]{Case 2} Suppose that $u_{2n-4}^{-1}(m)\not\in\{\overline n,n,m\}$ and that $\overline m, u_{2n-4}^{-1}(m),m$ appear in this (cyclic) order when we read the numbers on the circle in clockwise order. Let $Z$ be the set of numbers that either appear between $\overline m$ and $u_{2n-4}^{-1}(m)$ (including $\overline m$ and $u_{2n-4}^{-1}(m)$) or appear between $m$ and $\overline{u_{2n-4}^{-1}(m)}$ (including $m$ and $\overline{u_{2n-4}^{-1}(m)}$) in the clockwise order on the circle. Using \Cref{lem:DMainLemma}, it is straightforward to check that $u_{2n-4}$ cyclically permutes the elements of $Z$ in clockwise order and fixes each element of $\pm[n-1]\setminus Z$. Since $u_{2n-4}\in D_n$, we must have $u_{2n-4}(n)=\overline n$ and $u_{2n-4}(\overline n)=n$. Therefore, $u_{2n-4}=g(\rho)$, where $\rho$ is the type-$D_n$ noncrossing set partition whose blocks are $Z\cup\{\overline n,n\}$ and all of the singleton subsets of $\pm[n-1]\setminus Z$. Thus, $u_{2n-4}\in\NC(D_n,c)$. 
\end{enonce*}

\begin{enonce*}[remark]{Case 3} Suppose that $u_{2n-4}^{-1}(m)\not\in\{\overline n,n,m\}$ and that $m, u_{2n-4}^{-1}(m),\overline m$ appear in this (cyclic) order when we read the numbers on the circle in clockwise order. Let $B$ be the set of numbers that appear between $m$ and $u_{2n-4}^{-1}(m)$ (including $m$ and $u_{2n-4}^{-1}(m)$) in the clockwise order on the circle. Using \Cref{lem:DMainLemma}, it is straightforward to check that $u_{2n-4}$ cyclically permutes the elements of $B$ in clockwise order, cyclically permutes the elements of $\overline B$ in clockwise order, and fixes each element of $\pm[n-1]\setminus (B\cup(\overline B))$. Since $u_{2n-4}$ is in $ D_n$, it must fix $n$ and $\overline n$. Therefore, $u_{2n-4}=g(\rho)$, where $\rho$ is the type-$D_n$ noncrossing set partition whose blocks are $B$, $\overline B$, and all of the singleton subsets of $\pm[n]\setminus (B\cup(\overline B))$. Thus, $u_{2n-4}\in\NC(D_n,c)$.
\end{enonce*}

\begin{enonce*}[remark]{Case 4} Suppose that $u_{2n-4}^{-1}(m)\in\{\overline n,n\}$ and either that $u_{2n-4}^{-2}(m)=m$ or that the numbers $m,u_{2n-4}^{-2}(m),\overline m$ are distinct and appear in this (cyclic) order when we read the numbers on the circle in clockwise order. Let $B$ be the set of numbers that appear between $m$ and $u_{2n-4}^{-2}(m)$ (including $m$ and $u_{2n-4}^{-1}(m)$) in the clockwise order on the circle. Using \Cref{lem:DMainLemma}, it is straightforward to check that $u_{2n-4}$ cyclically permutes the elements of $B\cup\{u_{2n-4}^{-1}(m)\}$ in clockwise order, cyclically permutes the elements of $\overline B\cup\{\overline u_{2n-4}^{-1}(m)\}$ in clockwise order, and fixes each element of $\pm[n-1]\setminus (B\cup(\overline B))$. Therefore, $u_{2n-4}=g(\rho)$, where $\rho$ is the type-$D_n$ noncrossing set partition whose blocks are $B\cup\{u_{2n-4}^{-1}(m)\}$, $\overline B\cup\{\overline{u_{2n-4}^{-1}(m)}\}$, and all of the singleton subsets of $\pm[n-1]\setminus (B\cup(\overline B))$. Thus, $u_{2n-4}\in\NC(D_n,c)$. 
\end{enonce*}


\begin{enonce*}[remark]{Case 5} Suppose that $u_{2n-4}^{-1}(m)\in\{\overline n,n\}$ and either that $u_{2n-4}^{-2}(m)=\overline m$ or that the numbers $\overline m,u_{2n-4}^{-2}(m),m$ are distinct and appear in this (cyclic) order when we read the numbers on the circle in clockwise order. Let $r_1,r_2,\ldots,r_q$ be the numbers that appear between $\overline m$ and $u_{2n-4}^{-2}(m)$ (including $\overline m$ and $u_{2n-4}^{-2}(m)$) in the clockwise order on the circle. In particular, $r_1=\overline m$ and $r_q=u_{2n-4}^{-2}(m)$. In this case, one can use \Cref{lem:DMainLemma} to show that $u_{2n-4}$ is the single cycle 
\[
(r_1\,\,r_2\cdots r_q\,\, u_{2n-4}^{-1}(m)\,\,\overline {r_1}\,\,\overline {r_2}\cdots \overline{r_q}\,\,\overline{u_{2n-4}^{-1}(m)})
\]
 (meaning $u_{2n-4}$ fixes every element of $\pm[n]$ that is not in this cycle). However, this means that $u_{2n-4}$ has a single balanced cycle, contradicting the fact that every element of $D_n$ has an even number of balanced cycles. Therefore, this case cannot occur.  \qedhere
\end{enonce*}
\end{proof}

\begin{thebibliography}{99}
\bibitem[3]{armstrong2009generalized} Drew Armstrong, \emph{Generalized noncrossing partitions and combinatorics of Coxeter groups}, Mem. Amer. Math. Soc.202, no.949, American Mathematical Society, Providence, RI,2009.
\bibitem[6]{AthanasiadisReiner} Christos A. Athanasiadis and Victor Reiner, ``Noncrossing partitions for the group $D_n$'', \emph{SIAM J. Discrete Math.}\textbf{18}(2004),no.2,397--417.
\bibitem[7]{bessis2003dual} David Bessis, ``The dual braid monoid'', \emph{Ann. Sci. \'Ecole Norm. Sup.}(4)\textbf{36}(2003),no.5,647--683.
\bibitem[8]{bessis2015finite} David Bessis, ``Finite complex reflection arrangements are $K(\pi,1)$'', \emph{Ann. of Math.}(2)\textbf{181}(2015),no.3,809--904.
\bibitem[9]{BjornerBrenti} Anders Bj\"orner and Francesco Brenti, \emph{Combinatorics of Coxeter groups}, Graduate Texts in Mathematics231, Springer, New York,2005.
\bibitem[11]{BradyWatt} Thomas Brady and Colum Watt, ``$K(\pi,1)$'s for Artin groups of finite type'', \emph{Geom. Dedicata}\textbf{94}(2002),225--250.
\bibitem[12]{brady2008non} Thomas Brady and Colum Watt, ``Non-crossing partition lattices in finite real reflection groups'', \emph{Trans. Amer. Math. Soc.}\textbf{360}(2008),no.4,1983--2005.
\bibitem[16]{Carter} Roger W. Carter, ``Conjugacy classes in the Weyl group'', \emph{Compositio Math.}\textbf{25}(1972),1--59.
\bibitem[23]{ReadingNoncrossing} Nathan Reading, ``Noncrossing partitions, clusters and the Coxeter plane'', \emph{S\'em. Lothar. Combin.}\textbf{63}(2010),PaperB63b(32pages).
\bibitem[24]{reiner2017non} Victor Reiner, Vivien Ripoll and Christian Stump, ``On non-conjugate Coxeter elements in well-generated reflection groups'', \emph{Math. Z.}\textbf{285}(2017),no.3-4,1041--1062.
\end{thebibliography}
\end{document}
